\section{Objetivo general}
Aplicar técnicas de inteligencia artificial con enfoque híbrido para la identificación y categorización de segmentos de interés en señales bioeléctricas relacionadas con la epilepsia.

\section{Objetivos específicos}

\begin{itemize}

   \item Identificar las características más significativas de las señales bioeléctricas en los trabajos realizados previamente en UVG para detectar aquellas de mayor utilidad en este estudio.

    \item Implementar modelos con enfoque híbrido comunes para identificar la alternativa más robusta para el análisis de señales bioeléctricas relacionadas con la epilepsia.
    
    \item Evaluar el rendimiento del modelo híbrido mediante métricas de clasificación estándar, comparando su eficiencia con los resultados obtenidos en investigaciones previas en UVG.
    
    \item Actualizar la herramienta de software para el estudio de la epilepsia, desarrollada en fases anteriores, incorporando las mejoras al modelo para la detección de segmentos de interés de las señales bioeléctricas.

    
\end{itemize}

%\noindent
%Los objetivos se dividen en general y específicos.

%\noindent
%Puede aprender a redactar objetivos SMART en el siguiente \href{https://www.youtube.com/watch?v=DQlp0DvPBwU}{\underline{\textbf{enlace}}}.

%\noindent
%En caso no sepa qué verbos utilizar, consulte la taxonomía del siguiente \href{https://www.itson.mx/servicios/innovacion/documents/taxonomia_verbos_2.pdf}{\underline{\textbf{enlace}}}.



% Contenido original
% \section{Objetivo general}
% Los objetivos son las metas que se deseas alcanzar. Estos deben ser específicos, realistas, medibles y alcanzables. Además, inician con un verbo infinitivo, por ejemplo...

% \section{Objetivos específicos}

% \begin{itemize}
%     \item Implementar un sistema de purificación de glicerina, obtenida de la producción de biodiésel en el Laboratorio de Operaciones Unitarias.
    
%     \item Analizar la construcción del complejo de Edipo en \textit{Cien años de soledad}.  
    
%     \item Evaluar la adsorción y desorción de plomo (Pb) utilizando nanopartículas magnetizadas de carboximetilquitosano, sintetizadas a partir de desechos de camarón de la especie \textit{Litopenaeus vannamei}.
    
% \end{itemize}

% \noindent
% Los objetivos se dividen en general y específicos.

% \noindent
% Puede aprender a redactar objetivos SMART en el siguiente \href{https://www.youtube.com/watch?v=DQlp0DvPBwU}{\underline{\textbf{enlace}}}.

% \noindent
% En caso no sepa qué verbos utilizar, consulte la taxonomía del siguiente \href{https://www.itson.mx/servicios/innovacion/documents/taxonomia_verbos_2.pdf}{\underline{\textbf{enlace}}}.


